% !TEX spellcheck = en_US
% !TEX spellcheck = LaTeX
\documentclass[letterpaper,10pt,english]{article}
\usepackage{%
amsfonts,%
amsmath,%
amssymb,%
amsthm,%
babel,%
bbm,%
%biblatex,%
caption,%
centernot,%
color,%
enumerate,%
%enumitem,%
epsfig,%
epstopdf,%
etex,%
fancybox,%
framed,%
fullpage,%
%geometry,%
graphicx,%
hyperref,%
latexsym,%
mathptmx,%
mathtools,%
multicol,%
paralist,%
pgf,%
pgfplots,%
pgfplotstable,%
pgfpages,%
proof,%
psfrag,%
%subfigure,%
tcolorbox,%
tfrupee,%
tikz,%
times,%
%ulem,%
url,%
xcolor,%
mathpazo,%
}

\definecolor{shadecolor}{gray}{.95}%{rgb}{1,0,0}
\usepackage[margin=1in,top=0.75in]{geometry}
\usepackage[mathscr]{eucal}
\usepgflibrary{shapes}
\usepgfplotslibrary{fillbetween}
\usetikzlibrary{%
arrows,%
backgrounds,%
chains,%
decorations.pathmorphing,% /pgf/decoration/random steps | erste Graphik
decorations.text,% 
matrix,%
positioning,% wg. " of "
fit,%
patterns,%
petri,%
plotmarks,%
scopes,%
shadows,%
shapes.misc,% wg. rounded rectangle
shapes.arrows,%
shapes.callouts,%
shapes%
}

%\pgfplotsset{compat=newest} %<------ Here
\pgfplotsset{compat=1.11} %<------ Or use this one

\theoremstyle{plain}
\newtheorem{thm}{Theorem}[section]
\newtheorem{lem}[thm]{Lemma}
\newtheorem{prop}[thm]{Proposition}
\newtheorem{cor}[thm]{Corollary}
\newtheorem{clm}[thm]{Claim}

\theoremstyle{definition}
\newtheorem{defn}[thm]{Definition}
\newtheorem{conj}[thm]{Conjecture}
\newtheorem{exmp}[thm]{Example}
\newtheorem{axiom}[thm]{Axiom}
\newtheorem{assum}[thm]{Assumption}
\newtheorem{exerc}[thm]{Exercise}
\newtheorem*{defin}{Definition}

\theoremstyle{remark}
\newtheorem{rem}{Remark}
\newtheorem{note}{Note}

\newcommand{\iid}{\emph{i.i.d.}\ }
\newcommand{\Cov}{\operatorname{Cov}}
%\newcommand{\det}{\operatorname{det}}
%\newcommand{\Real}{\mathbb{R}}
\newcommand{\tr}{\operatorname{tr}}
\newcommand{\Var}{\operatorname{Var}}
\newcommand{\cov}{\operatorname{cov}}

%\renewcommand{\proof}[1]{\begin{proof}#1\end{proof}}
\newcommand{\EQ}[1]{\begin{equation*}#1\end{equation*}}
\newcommand{\EQN}[1]{\begin{equation}#1\end{equation}}
\newcommand{\eq}[1]{\begin{align*}#1\end{align*}}
\newcommand{\meq}[2]{\begin{xalignat*}{#1}#2\end{xalignat*}}
\newcommand{\norm}[1]{\left\lVert#1\right\rVert}
\newcommand{\inner}[1]{\left\langle #1\right\rangle}
\newcommand{\abs}[1]{\left\lvert#1\right\rvert}
\newcommand{\expect}[1]{\mathbb{E}\left[{#1}\right]}
\newcommand{\prob}[1]{\mathbb{P}\left[{#1}\right]}
\newcommand{\given}{\; \big\vert \;} 
\newcommand{\set}[1]{\left\{#1\right\}} 
\newcommand{\indicator}[1]{1_{\set{#1}}} 
\newcommand{\Ind}[1]{\mathbbm{1}_{#1}} 
\newcommand{\SetIn}[1]{\mathbbm{1}_{\set{#1}}} 
\newcommand{\red}[1]{\textcolor{red}{#1}} 
\newcommand{\floor}[1]{\left\lfloor#1\right\rfloor}
\newcommand{\ceil}[1]{\left\lceil#1\right\rceil}


\newcommand{\C}{\mathbb{C}}
\newcommand{\D}{\mathbb{D}}
\newcommand{\E}{\mathbb{E}}
\newcommand{\F}{\mathbb{F}}
\newcommand{\N}{\mathbb{N}}
\renewcommand{\P}{\mathbb{P}}
\newcommand{\Q}{\mathbb{Q}}
\newcommand{\R}{\mathbb{R}}
\newcommand{\Z}{\mathbb{Z}}

\newcommand{\bA}{\mathbf{A}}
\newcommand{\bI}{\mathbf{I}}
\newcommand{\bU}{\mathbf{1}}
\newcommand{\bx}{\mathbf{x}}
\newcommand{\bX}{\mathbf{X}}
\newcommand{\bZ}{\mathbf{Z}}


\newcommand{\cA}{\mathcal{A}}
\newcommand{\cB}{\mathcal{B}}
\newcommand{\cC}{\mathcal{C}}
\newcommand{\cF}{\mathcal{F}}
\newcommand{\cG}{\mathcal{G}}
\newcommand{\cM}{\mathcal{M}}
\newcommand{\cN}{\mathcal{N}}
\newcommand{\cP}{\mathcal{P}}
\newcommand{\cT}{\mathcal{T}}
\newcommand{\cX}{\mathcal{X}}
\newcommand{\cY}{\mathcal{Y}}

\newcommand{\sA}{\mathscr{A}}
\newcommand{\sB}{\mathscr{B}}
\newcommand{\sC}{\mathscr{C}}
\newcommand{\sE}{\mathscr{E}}
\newcommand{\sF}{\mathscr{F}}
\newcommand{\sG}{\mathscr{G}}
\newcommand{\sH}{\mathscr{H}}
\newcommand{\sL}{\mathscr{L}}
\newcommand{\sS}{\mathscr{S}}
\newcommand{\sT}{\mathscr{T}}
\newcommand{\sX}{\mathscr{X}}
\newcommand{\sY}{\mathscr{Y}}
\newcommand{\sZ}{\mathscr{Z}}

\newcommand{\tS}{\tilde{S}}
% Debug
\newcommand{\todo}[1]{\begin{color}{blue}{{\bf~[TODO:~#1]}}\end{color}}

% a few handy macros

\renewcommand{\le}{\leqslant}
\renewcommand{\ge}{\geqslant}
\newcommand\matlab{{\sc matlab}}
\newcommand{\goto}{\rightarrow}
\newcommand{\bigo}{{\mathcal O}}
%\newcommand{\half}{\frac{1}{2}}
%\newcommand\implies{\quad\Longrightarrow\quad}
\newcommand\reals{{{\rm l} \kern -.15em {\rm R} }}
\newcommand\complex{{\raisebox{.043ex}{\rule{0.07em}{1.56ex}} \hskip -.35em {\rm C}}}


% macros for matrices/vectors:

% matrix environment for vectors or matrices where elements are centered
\newenvironment{mat}{\left[\begin{array}{ccccccccccccccc}}{\end{array}\right]}
\newcommand\bcm{\begin{mat}}
\newcommand\ecm{\end{mat}}

% matrix environment for vectors or matrices where elements are right justifvied
\newenvironment{rmat}{\left[\begin{array}{rrrrrrrrrrrrr}}{\end{array}\right]}
\newcommand\brm{\begin{rmat}}
\newcommand\erm{\end{rmat}}

% for left brace and a set of choices
%\newenvironment{choices}{\left\{ \begin{array}{ll}}{\end{array}\right.}
\newcommand\when{&\text{if~}}
\newcommand\otherwise{&\text{otherwise}}
% sample usage:
%\delta_{ij} = \begin{choices} 1 \when i=j, \\ 0 \otherwise \end{choices}


% for labeling and referencing equations:
\newcommand{\eql}{\begin{equation}\label}
\newcommand{\eqn}[1]{(\ref{#1})}
% can then do
%\eql{eqnlabel}
%...
%\end{equation}
% and refer to it as equation \eqn{eqnlabel}.
% some useful macros for finite difference methods:
\newcommand\unp{U^{n+1}}
\newcommand\unm{U^{n-1}}
% for chemical reactions:
\newcommand{\react}[1]{\stackrel{K_{#1}}{\rightarrow}}
\newcommand{\reactb}[2]{\stackrel{K_{#1}}{~\stackrel{\rightleftharpoons}
 {\scriptstyle K_{#2}}}~}
\makeatletter
\def\th@plain{%
\thm@notefont{}% same as heading font
\itshape % body font
}
\def\th@definition{%
\thm@notefont{}% same as heading font
\normalfont % body font
}
\makeatother
\date{}


\title{Lecture-04: Random Variable}
\author{}

\begin{document}
\maketitle

\section{Random Variable}
%\begin{defn}[Borel sets] 
%The smallest $\sigma$-algebra generated by the half-open sets $B_x = (-\infty, x]$ for all $x \in \R$ is called a \textbf{Borel $\sigma$-algebra} and denoted by $\sB(\R) = \sigma(\set{B_x: x \in \R})$. 
%The elements of the Borel $\sigma$-algebra are called \textbf{Borel sets}. 
%\end{defn}
\begin{defn}[Random variable]
Consider a probability space $(\Omega, \sF, P)$. 
A \textbf{random variable} $X: \Omega \to \R$ is a real-valued function from the sample space to real numbers, 
such that for each $x \in \R$ %the event space $\sF$ is generated by 
the event 
\EQ{
A_X(x) \triangleq \set{\omega \in \Omega: X(\omega) \le x}
= \set{X \le x} = X^{-1}(-\infty, x] = X^{-1}(B_x) \in \sF.
}  
%In particular, we consider a probability space $(\Omega, \sF, P)$ where $E_x \in \sF$ for each real $x \in \R$. 
We say that the random variable $X$ is $\sF$-measurable.  
\end{defn}
\begin{rem}
Recall that the set $A_X(x)$ is always a subset of sample space $\Omega$ for any mapping $X:\Omega\to\R$, 
and $A_X(x)\in\sF$ is an event when $X$ is a random variable.
\end{rem}
\begin{shaded}
\begin{exmp}[Constant function] 
Consider a mapping $X:\Omega\to \set{c} \subseteq \R$ defined on an arbitrary probability space $(\Omega, \sF, P)$, 
such that $X(\omega) = c$ for all outcomes $\omega \in \Omega$. 
We observe that 
\EQ{
A_X(x) = X^{-1}(B_x) = 
\begin{cases}
\emptyset, & x < c,\\
\Omega, & x \ge c.
\end{cases}
}
That is $A_X(x) \in \sF$ for all event spaces, and hence $X$ is a random variable and measurable for all event spaces. 
\end{exmp}
\begin{exmp}[Indicator function] 
For an arbitrary probability space $(\Omega,\sF, P)$ and an event $A \in \sF$, 
consider the indicator function $\Ind{A}:\Omega \to [0,1]$. 
Let $x \in \R$, and $B_x = (-\infty, x]$, then it follows that 
\EQ{
A_X(x) = \Ind{A}^{-1}(B_x) 
= \begin{cases}
\Omega, & x \ge 1,\\
A^c, & x \in [0, 1),\\
\emptyset, & x < 0.
\end{cases} 
}  
That is, $A_X(x) \in \sF$ for all $x \in \R$, and hence the indicator function $\Ind{A}$ is a random variable.  
\end{exmp}
\end{shaded}

\begin{rem} 
Since any outcome $\omega \in \Omega$ is random, so is the real value $X(\omega)$. % \in \R$. 
\end{rem}
\begin{rem}
Probability is defined only for events and not for random variables. 
The events of interest for random variables are the \textbf{lower level sets} $A_X(x) = \set{\omega: X(\omega) \le x} = X^{-1}(B_x)$ for any real $x$. % \in \R$. 
\end{rem}
\begin{rem}
Consider a probability space $(\Omega,\sF,P)$ and a random variable $X:\Omega \to \R$ that is $\sG$ measurable for $\sG \subseteq \sF$. 
If $\sG \subseteq \sH$, then $X$ is also $\sH$ measurable.
\end{rem}
%\begin{defn}[Indicator functions] 
%For a probability space $(\Omega, \sF, P)$ and any event $A \in \sF$, 
%we can define an \textbf{indicator function} $\Ind{A}: \Omega \to \set{0,1}$ as 
%\EQ{
%\Ind{A}(\omega) = \begin{cases}
%1, & \omega \in A,\\
%0, & \omega \notin A.
%\end{cases}
%} 
%\end{defn}
%\begin{shaded}
%\begin{exmp} 
%We will show that indicator function $\Ind{A}$ is a random variable for any event $A \in \sF$.  
%Let $x \in \R$, and $B_x = (-\infty, x]$, then it follows that 
%\EQ{
%\Ind{A}^{-1}(B_x) 
%= \begin{cases}
%\Omega, & x \ge 1,\\
%A^c, & x \in [0, 1),\\
%\emptyset, & x < 0.
%\end{cases} 
%}  
%That is, $\Ind{A}^{-1}(B_x) \in \sF$ for all $x \in \R$, and hence it is a random variable. 
%Further, we can write $\sigma(\Ind{A}) = \set{\emptyset, A, A^c, \Omega}$. 
%In addition, we have 
%\EQ{
%F_X(x) = 
%\begin{cases}
%1, & x \ge 1,\\
%1- P(A), & x \in [0, 1),\\
%0, & x < 0.
%\end{cases}
%}
%\end{exmp}
%\end{shaded}
\subsection{Distribution function for a random variable}
\begin{defn}
For an $\sF$ measurable random variable $X:\Omega\to\R$ defined on the probability space $(\Omega,\sF,P)$, 
we can associate a \textbf{distribution function} (CDF) $F_X: \R \to [0,1]$ 
such that for all $x \in \R$, 
\EQ{
F_X(x) \triangleq P(A_X(x)) = P(\set{X \le x}) = P \circ X^{-1}(-\infty, x] = P\circ X^{-1}(B_x).
}
%, and specifies the probability induced by the random variable.
\end{defn}
\begin{shaded}
\begin{exmp}[Constant random variable] 
Let $X:\Omega \to \set{c} \subseteq \R$ be a constant random variable defined on the probability space $(\Omega, \sF, P)$. 
The distribution function is a right-continuous step function at $c$ with step-value unity. 
That is, $F_X(x) = \Ind{[c,\infty)}(x)$.  
We observe that $P(\set{X = c}) = 1$. 
%Does this make sense? Is $\set{\omega \in \Omega: X(\omega) = x}$ an event for a general random variable? 
%
%Consider a monotonically increasing sequence $x\in\R^\N$ such that $\lim_nx_n = x_0$. 
%Then, the sequence of events $(A_{x_n} \in \sF: n \in \N)$ is monotonically increasing and hence the $\lim_nA_{x_n} = \set{X < x_0} \in \sF$ is an event. 
%It follows that $\set{X = x_0} = A_{x_0}\cap\set{X < x_0}^c \in \sF$ is also an event. 
\end{exmp}
%\end{shaded}
%\begin{shaded}
\begin{exmp}[Indicator random variable] 
For an indicator random variable $\Ind{A}:\Omega \to \set{0,1}$ defined on a probability space $(\Omega, \sF, P)$ and an event $A \in \sF$,  
we have 
\EQ{
F_X(x) = 
\begin{cases}
1, & x \ge 1,\\
1- P(A), & x \in [0, 1),\\
0, & x < 0.
\end{cases}
}
\end{exmp}
\end{shaded}
\begin{lem}[Properties of distribution function] 
The distribution function $F_X$ for any random variable $X$ satisfies the following properties. 
\begin{enumerate}
\item The distribution function is monotonically non-decreasing in $x \in \R$. 
\item The distribution function is right-continuous at all points $x \in \R$.
\item The upper limit is $\lim_{x \to \infty}F_X(x) = 1$ and the lower limit is $\lim_{x \to -\infty}F_X(x) = 0$. 
\end{enumerate}
\end{lem}
\begin{proof}
Let $X$ be a random variable defined on the probability space $(\Omega, \sF, P)$. 
\begin{enumerate}
\item Let $x_1, x_2 \in \R$ such that $x_1 \le x_2$.  
Then for any $\omega \in A_{x_1}$, we have $X(\omega) \le x_1 \le x_2$, and it follows that $\omega \in A_{x_2}$. 
This implies that $A_{x_1} \subseteq A_{x_2}$. 
The result follows from the monotonicity of the probability. 
 
\item For any $x \in \R$, consider any monotonically decreasing sequence $x\in\R^\N$ such that $\lim_n x_n = x_0$. 
It follows that the sequence of events $\big(A_{x_n} = X^{-1}(-\infty, x_n] \in \sF: n \in \N\big)$, is monotonically decreasing and hence $\lim_{n \in \N}A_{x_n} = \cap_{n \in \N}A_{x_n} = A_{x_0}$. 
The right-continuity then follows from the continuity of probability, since 
\EQ{
F_X(x_0) = P(A_{x_0}) = P(\lim_{n \in \N}A_{x_n}) = \lim_{n \in \N}P(A_{x_n}) = \lim_{x_n \downarrow x}F(x_n). 
}
\item Consider a monotonically increasing sequence $x\in\R^\N$ such that $\lim_{n}x_n = \infty$, then $(A_{x_n} \in \sF: n \in \N)$ is a monotonically increasing sequence of sets and $\lim_nA_{x_n} = \cup_{n \in \N}A_{x_n} = \Omega$. 
From the continuity of probability, it follows that 
\EQ{
\lim_{x_n \to \infty}F_X(x_n) = \lim_nP(A_{x_n}) = P(\lim_n A_{x_n}) = P(\Omega) = 1.
}
Similarly, we can take a monotonically decreasing sequence $x\in \R^\N$ such that $\lim_{n}x_n = -\infty$, 
then $(A_{x_n} \in \sF: n \in \N)$ is a monotonically decreasing sequence of sets and $\lim_nA_{x_n} = \cap_{n \in \N}A_{x_n} = \emptyset$. 
From the continuity of probability, it follows that $\lim_{x_n \to -\infty}F_X(x_n) = 0$. 
\end{enumerate}
\end{proof}
\begin{rem}
If two reals $x_1 < x_2$ then $F_X(x_1) \le F_X(x_2)$ with equality if and only if $P\set{(x_1 < X \le x_2}) = 0$. 
This follows from the fact that $A_{x_2} = A_{x_1}\cup X^{-1}(x_1, x_2]$. 
\end{rem}
\subsection{Event space generated by a random variable}
\begin{defn}[Event space generated by a random variable] 
Let $X: \Omega \to \R$ be an $\sF$ measurable random variable defined on the probability space $(\Omega, \sF, P)$. 
The smallest event space generated by the events $A_X(x) = X^{-1}(B_x) = X^{-1}(-\infty, x]$ for $x \in \R$ is called the \textbf{event space generated} by this random variable $X$, 
and denoted by $\sigma(X) \triangleq \sigma(\set{A_X(x): x \in \R})$. 
\end{defn}
\begin{rem}
The event space generated by a random variable is the collection of the inverse of Borel sets, i.e. $\sigma(X) = \set{X^{-1}(B): B \in \sB(\R)}$.
This follows from the fact that $A_X(x) = X^{-1}(B_x)$ and the inverse map respects countable set operations such as unions, complements, and intersections. 
That is, if $B \in \sB(\R) = \sigma(\set{B_x:x\in \R})$, then $X^{-1}(B) \in \sigma(\set{A_X(x): x \in \R})$. 
Similarly, if $A \in \sigma(X) = \sigma(\set{A_X(x): x \in \R})$, then $A = X^{-1}(B)$ for some $B \in \sigma(\set{B_x: x \in \R})$.  
\end{rem}
\begin{shaded}
\begin{exmp}[Constant random variable] 
Let $X:\Omega \to \set{c} \subseteq \R$ be a constant random variable defined on the probability space $(\Omega,\sF,P)$.  
Then the smallest event space generated by this random variable is 
$\sigma(X) = \set{\emptyset, \Omega}$. 
\end{exmp}
\begin{exmp}[Indicator random variable] 
Let $\Ind{A}$ be an indicator random variable defined on the probability space $(\Omega,\sF,P)$ and event $A \in \sF$, 
then the smallest event space generated by this random variable is 
$\sigma(X) = \sigma(\set{\emptyset, A^c, \Omega}) = \set{\emptyset, A, A^c, \Omega}$. 
\end{exmp}
\end{shaded}



\subsection{Discrete random variables}
\begin{defn}[Discrete random variables] 
If a random variable $X: \Omega \to \sX \subseteq \R$ takes countable values on real-line, 
then it is called a \textbf{discrete random variable}. 
That is, the range of random variable $\sX$ is countable, 
and the random variable is completely specified by the \textbf{probability mass function} 
\EQ{
P_X(x) = P(\set{X = x}), \text{ for all }x \in \sX. 
}
\end{defn}
\begin{shaded}
\begin{exmp}[Bernoulli random variable] 
For the probability space $(\Omega, \sF, P)$, the \textbf{Bernoulli random variable} is a mapping $X: \Omega \to \set{0,1}$ and $P_X(1) = p$. 
We observe that Bernoulli random variable is an indicator for the event $A \triangleq X^{-1}\set{1}$, and $P(A) = p$. 
Therefore,  the distribution function $F_X$ is given by 
\EQ{
F_X = (1-p)\Ind{[0, 1)} + \Ind{[1, \infty)}.
}
\end{exmp}
\end{shaded}
\begin{lem} 
Any discrete random variable is a linear combination of indicator function over a partition of the sample space.
\end{lem}
\begin{proof}
For a discrete random variable $X: \Omega \to \sX \subset \R$ on a probability space $(\Omega, \sF, P)$, 
the range $\sX$ is countable, 
and we can define events $E_x \triangleq \set{\omega \in \Omega: X(\omega) = x} \in \sF$ for each $x \in \sX$. 
Then the mutually disjoint sequence of events $(E_x \in \sF: x \in \sX)$ partitions the sample space $\Omega$.  %and are events from the definition of random variables. 
We can write 
\EQ{
X(\omega) = \sum_{x \in \sX}x\Ind{E_x}(\omega).
}
\end{proof}
\begin{defn} 
Any discrete random variable $X:\Omega \to \sX \subseteq\R$ defined over a probability space $(\Omega,\sF,P)$, 
with finite range is called a simple random variable. 
\end{defn}
\begin{shaded}
\begin{exmp}[Simple random variables] 
Let $X$ be a simple random variable, then $X = \sum_{x \in \sX}x\Ind{A_X(x)}$ where $(A_X(x) = X^{-1}\set{x} \in \sF: x \in \sX)$ is a finite partition of the sample space $\Omega$. 
Without loss of generality, we can denote $\sX  = \set{x_1, \dots, x_n}$ where $x_1 \le \dots \le x_n$. 
Then, 
\EQ{
X^{-1}(-\infty, x] = \begin{cases}
\Omega, &x \ge x_n,\\
\cup_{j=1}^iA_X(x_j), & x \in [x_i, x_{i+1}), i \in [n-1],\\
\emptyset, &x < x_1.
\end{cases}
}
Then the smallest event space generated by the simple random variable $X$ is $\set{\cup_{x \in S}A_X(x): S \subseteq \sX}$. 
\end{exmp}
\end{shaded}
%\begin{shaded}
%\begin{exmp} 
%We will show that indicator function $\Ind{A}$ is a discrete random variable for any event $A \in \sF$.   
%\end{exmp}
%\end{shaded}
%
%\begin{defn}[Simple functions]
%For a probability space $(\Omega, \sF, P)$, a finite $n$, events $E_1, \dots, E_n \in \sF$, and real constants $x_1, \dots, x_n$
%we can define a simple function $X: \Omega \to \set{x_1, \dots, x_n}$ to be 
%\EQ{
%X(\omega) = \sum_{i=1}^nx_i\Ind{E_i}(\omega). 
%}
%\end{defn}

\subsection{Continuous random variables}
\begin{defn}
For a continuous random variable $X$, 
there exists \textbf{density function} $f_X: \R \to [0, \infty)$ such that 
\EQ{
F_X(x) = \int_{-\infty}^xf_X(u)du.
}
\end{defn}
\begin{shaded}
\begin{exmp}[Gaussian random variable]
For a probability space $(\Omega, \sF, P)$, \textbf{Gaussian random variable} is a continuous random variable $X: \Omega \to \R$ defined by its density function 
\EQ{
f_X(x) = \frac{1}{\sqrt{2\pi}\sigma}\exp\left(-\frac{(x-\mu)^2}{2\sigma^2}\right),~ x \in \R.
}
\end{exmp}
\end{shaded}

\end{document}
